\documentclass[10pt,a4paper]{amsart}
\usepackage[margin=19mm]{geometry}
\usepackage{amsmath,amssymb,mathtools}
\usepackage[scr=rsfs,cal=euler]{mathalfa}
\usepackage{xfrac}
\usepackage[unicode,hidelinks,bookmarks=false]{hyperref}
\newcommand{\ie}{i.e.\ }
\newcommand{\cf}{cf.\ }
\newcommand{\resp}{respectively\ }
\newcommand{\Subsection}{\subsection}
\providecommand{\nobreakdash}{\nobreak\hbox{-}}
\setlength{\emergencystretch}{2em}
\multlinegap0pt
\usepackage{amssymb}
\newtheorem{proposition}{Proposition}
\newcommand{\C}{\mathbb{C}}
\newcommand{\Res}{\operatorname{Res}}
\newcommand{\Z}{\mathbb{Z}}
\newcommand{\bfb}{{\boldsymbol{b}}}
\newcommand{\bff}{{\boldsymbol{f}}}
\newcommand{\bfxi}{{\boldsymbol{\xi}}}
\newcommand{\cA}{{\mathcal A}}
\newcommand{\cD}{{\mathcal D}}
\newcommand{\cO}{{\mathcal O}}
\title{Toric Euler-Jacobi vanishing theorem and zeros at infinity}
\author{Carlos D'Andrea, Alicia Dickenstein}
\date{Source: arXiv:2601.13977v2 (CC BY 4.0)}
\begin{document}
\maketitle

\noindent\emph{Source and licence.} Toric Euler-Jacobi vanishing theorem and zeros at infinity. Version arXiv:2601.13977v2 (CC BY 4.0), distributed by arXiv under the Creative Commons Attribution 4.0 International licence, http://creativecommons.org/licenses/by/4.0/, as recorded in the arXiv licence metadata for this exact version and shown on the abstract page https://arxiv.org/abs/2601.13977v2. Full licence notice: \texttt{licenses/arxiv-2601.13977-CC-BY-4.0.txt}.\par\smallskip

\noindent\textit{Editorial scope.} Counted unit: the article's proposition neq (source0.tex lines 1027--1029). The article's complete original proof of it is delivered with it (source0.tex lines 1030--1054, one \texttt{proof} environment of 25 lines). No mathematics is written, repaired or re-derived: the statement and the proof are the article's own lines, in the article's own notation, and no retained line is substituted or re-flowed.  Retained uncounted as the source-local context the proof needs: lines 810--813; lines 927--929; lines 1016--1018; lines 1022--1024.  Nothing was omitted from any retained span, so the package carries no omission record.  No retained line contains a citation, so the wrapper carries no bibliography and no bibliography entry.  No retained line contains a figure environment, an image-inclusion command or drawing code, so no image is delivered and the package declares no asset dependency.  Editorial formatting only, in addition: an AMS \texttt{amsart} preamble that restates the article's own declarations and macros used by the retained lines, namely the environments \texttt{proposition} on separate counters, together with the standard packages those lines need.  The retained environments print in this document as Proposition 1 in order of appearance, whereas the article prints the same items with the numbering of its own full document.  The complete native source of the article as distributed in the arXiv e-print for this exact version is bundled unchanged as \texttt{sources/193061/source0.tex}.
\begin{proposition}\label{lresid}
Let $h$ be such that $h_\bfxi\in\cO_\bfxi,$ Then,
 $h_\bfxi\in I_\bfxi\iff R_{\bff, \bfxi}\big[\frac{h\,dz}{f_1\ldots f_n}\big]\equiv0,$ the zero operator. 
\end{proposition}

\begin{equation}\label{lbtz}
\partial(g\cdot f)=\sum_{\bfb\in\Z_{\geq0}^\ell} \partial_\bfb(g)\big(s_\bfb(\partial)\big)(f) \ \forall \partial\in\cD,\, \forall g,\,f\in\C[y_1,\ldots, y_\ell].
\end{equation}

\begin{equation}\label{zizi}
\C[X_\sigma]/\langle F_1,\ldots, F_n\rangle\simeq \C[u_1,\ldots, u_n]^{G(\sigma)}/\langle \tilde{F}_1,\ldots, \tilde{F}_n\rangle\simeq \C[y_1,\ldots, y_\ell]/ I_{\sigma,\cA},
\end{equation}

\begin{equation}\label{aunt}
\mbox{Res}_{F,\bfxi_i}=\sum_{j=1}^{\mu_i} c_{i,j}\partial_{i,j}.
\end{equation}

\begin{proposition}\label{neq}
With notation as in \eqref{aunt}, we have that $c_{i,\mu_i}\neq0.$
\end{proposition}

\begin{proof}
Suppose that 
\begin{equation}\label{an}
\mbox{Res}_{F,\bfxi_i}=\sum_{j=1}^{\mu_i-1} c_{i,j}\partial_{i,j}.
\end{equation}
 As $\{\partial^{i,1},\ldots, \partial^{i,\mu_i}\}$ is a basis of $V_i,$ there exists $G\notin Q_i$ such that $\partial_{i,j}(G)=0$ if $j<\mu_i,$ and $\partial_{i,\mu_i}(G)\neq0.$   Via \eqref{zizi} we assume w.l.o.g. that $G\in\C[u_1,\ldots, u_n]^{G(\sigma)}.$ As $G\notin Q_i,$ from Proposition \ref{lresid} we deduce that  there is $H\in\C[u_1,\ldots, u_n]$ such that
\begin{equation}\label{deux}
\Res_{\bf0}\left(\frac{G \cdot H}{\tilde{F}_{1}\ldots \tilde{F}_n}\,d u_1\wedge\ldots\wedge d u_n\right)\neq0.
\end{equation}
As $H$ is not necessarily in $\C[u_1,\ldots, u_n]^{G(\sigma)},$ we will replace it with 
$$H^\sigma:=\sum_{g\in G(\sigma)} H(g(u))\, g(u_1)\ldots g(u_n)\in\C[u_1,\ldots, u_n]^{G(\sigma)},$$ and note that
$$\begin{array}{ccl}
\Res_{\bf0}\left(\frac{G \cdot H^\sigma}{\tilde{F}_{1}\ldots \tilde{F}_n}\frac{d u_1}{u_1}\wedge\ldots\wedge \frac{d u_n}{u_n}\right)& = &
\sum_{g\in G(\sigma)}\Res_{\bf0}\left(\frac{G \cdot H(g(u))}{\tilde{F}_{1}\ldots \tilde{F}_n}d g(u_1)\wedge\ldots\wedge dg( u_n)\right)\\
&=& |G(\sigma)|\Res_{\bf0}\left(\frac{G \cdot H}{\tilde{F}_{1}\ldots \tilde{F}_n}d u_1\wedge\ldots\wedge d u_n\right)\neq0,
\end{array}
$$
so we actually have that $\Res_{F,\bfxi_i}(G\,H^\sigma)\neq0.$ 

In addition, by construction we have $H^\sigma(\bf0)=0.$ 
This consideration, combined with Leibniz's rule \eqref{lbtz} applied to $\partial_{i,j}(G\cdot H^\sigma),$ and using that the subspace $V_i$ is closed,  and that $\partial_{i,j}(G)=0$ if $j<\mu_i,$  shows that 
 $\partial_{i,j}(G\cdot H^\sigma)=0$ for all $j<\mu_i.$ But this implies that 
 $$0\neq \Res_{F,\bfxi_i}(G\,H^\sigma)=\sum_{j=1}^{\mu_j-1}c_{i,j} \partial_{i, j}(G\,H^\sigma)=0,$$ 
a contradiction.
\end{proof}

\end{document}
